\documentclass[11pt]{article}
\usepackage[margin=1in]{geometry}
\usepackage{amsmath,amssymb,amsthm,mathtools}
\IfFileExists{lmodern.sty}{\usepackage{lmodern}}{}
\usepackage[expansion=false]{microtype}
\usepackage{booktabs,array}
\usepackage{enumitem}
\usepackage{xcolor}
\usepackage[colorlinks=true,linkcolor=blue!50!black,citecolor=green!35!black,urlcolor=blue!50!black]{hyperref}
\usepackage[nameinlink]{cleveref}

\newtheorem{theorem}{Theorem}[section]
\newtheorem{proposition}[theorem]{Proposition}
\newtheorem{lemma}[theorem]{Lemma}
\newtheorem{corollary}[theorem]{Corollary}
\theoremstyle{definition}
\newtheorem{definition}[theorem]{Definition}
\newtheorem{example}[theorem]{Example}
\theoremstyle{remark}
\newtheorem{remark}[theorem]{Remark}
\newtheorem*{mainthm}{Main results}

\newcommand{\D}{\mathsf{D}}
\newcommand{\R}{\mathbb{R}}
\newcommand{\Z}{\mathbb{Z}}
\newcommand{\Cpl}{\operatorname{Cpl}}
\newcommand{\diag}{\operatorname{diag}}
\newcommand{\T}{\mathsf T}
\newcommand{\Hess}{\operatorname{Hess}}

\title{Bridges:\\ From Probability and Entropy to Branes, with Worked Examples}
\author{Vinay K. Awasthi\\
Department of Applied Mathematics and Statistics\\
Johns Hopkins University\thanks{\texttt{check\_bridges.py} (9 checks) verifies numerically the computable statements proved here, on examples
and random instances, and every number in the worked examples. Propositions 4.2, 4.3 and the derived invariants in
Proposition 6.1 rest on the cited literature; results cited from the other papers are checked there.}}
\date{October 8, 2026, draft 1}

\begin{document}
\maketitle

\begin{abstract}
We construct the maps between five layers, probability, information geometry, entropy, couplings and symplectic branes,
and test which structures they preserve. Entropy's negative Hessian is the Fisher metric. Marginal constraints become
Lagrangian branes as conormal bundles, and the graph of the entropy differential meets such a brane exactly at the
maximum-entropy law of the constraint, transversally because the Fisher metric is nondegenerate. When no positive law meets
the constraint the brane is empty. In a worked example with three bits this happens exactly when the pairwise disagreement
reaches $2/3$, the contextual fraction is $1-3\varepsilon$ below the threshold, and the Fisher form of the fibre blows up as
the threshold is approached. Because a cell of the filler vanishes there, joint samples detect the boundary at distance $m$
from $O(1/m)$ samples, while pairwise-only samples need $O(1/m^2)$. On the toric side,
entropy is an admissible symplectic potential on the toric manifold of every nondegenerate coupling polytope, without the
facet hypothesis of earlier work, so couplings are the points of a positive real Lagrangian and the bases of torus branes. When the
toric manifold is Fano, only the torus branes at critical points of the disc potential are nonzero objects; for $2\times2$
couplings this is the Fréchet midpoint, which is generally not the product coupling. Gluing with a fixed coupling acts on couplings linearly, and it preserves the lattice of the
torus exactly when its kernel is deterministic or constant. So gluing with a non-constant stochastic kernel is not induced by a
torus-equivariant correspondence, a defect related to, though not identical with, the defects of stochastic kernels found in
other layers of the programme.
\end{abstract}

\tableofcontents

%=====================================================================
\section{Introduction}\label{sec:intro}
%=====================================================================

The programme of which this note is part works with five layers:
\begin{itemize}[leftmargin=*,nosep]
\item probability distributions, which describe states and uncertainty;
\item information geometry, which measures how distinguishable nearby states are;
\item entropy, which supplies a potential and selects canonical distributions;
\item coupling spaces, which describe compatible joint realizations \cite{Conditionals,Algebras};
\item symplectic branes and their Fukaya-type categories, which could organize admissible structures, transformations and
  their higher compositions \cite[\S\S7, 10]{Conditionals}.
\end{itemize}
The task is to construct the maps between the layers and to find out which structures they preserve. Each bridge below
comes with a proof and a worked example in which every number can be checked by hand; the one bridge where data enter, the
detection of a breaking point, comes with a synthetic test.

\begin{mainthm}
\leavevmode
\begin{enumerate}[label=(\Alph*),leftmargin=*]
\item \emph{Constraint branes} (\cref{thm:conormal}). The conormal bundle of a marginal constraint is a Lagrangian in the
  cotangent bundle of the open simplex. The graph of $dH$ meets it in exactly one point, the maximum-entropy law of the
  constraint, transversally; it misses it exactly when no positive law satisfies the constraint.
\item \emph{Three bits} (\cref{ex:threebits}). For pairwise disagreement $1-\varepsilon$, a joint exists iff $\varepsilon\ge\frac13$, a
  positive one iff $\frac13<\varepsilon<1$; the contextual fraction is $1-3\varepsilon$ below $\frac13$; the Fisher form of the fibre is
  $\frac8{3\varepsilon-1}+\frac{24}{1-\varepsilon}$; and from joint samples the threshold is detected at distance $m$ with $O(1/m)$ samples.
\item \emph{Toric branes} (\cref{thm:abreu,prop:realpositive,prop:floer}). On the toric manifold of every nondegenerate
  coupling polytope, entropy is an admissible symplectic potential; couplings are the points of the positive real Lagrangian
  and the bases of the torus branes; when the manifold is Fano, only torus branes at critical points of the disc potential
  survive, which for $2\times2$ couplings is the Fréchet midpoint (\cref{ex:twobytwo}).
\item \emph{Correspondences} (\cref{thm:integral,ex:gluetwo}). Gluing with a coupling whose kernel is $K$ maps lattice to
  lattice iff $K$ is deterministic or constant, so a non-constant stochastic $K$ is not induced by an equivariant correspondence; in the $2\times2$ case it scales the moment segment by the Dobrushin coefficient.
\item \emph{Categorical layer} (\cref{prop:morita}). The gluing algebra is Morita equivalent to $\R\times\R$, so all its derived
  invariants are those of two points.
\end{enumerate}
\end{mainthm}

\paragraph{What is not claimed.} We do not construct a functor from couplings to a Fukaya category. We show where the obvious
candidates fail and why: in the cotangent bundle of the simplex, which is contractible, the entropy brane carries only its one
intersection point; in the Fano case, torus branes over most couplings are zero objects; and gluing with a non-constant
stochastic kernel is not induced by an equivariant correspondence of the kind defined in \cref{sec:correspondences}.
Non-equivariant correspondences are not ruled out.

%=====================================================================
\section{From probability to geometry}\label{sec:geometry}
%=====================================================================

Let $\Delta^\circ$ be the open simplex of full-support distributions on $n$ points, with tangent space
$T=\{u:\sum u_y=0\}$, and $H(p)=-\sum p_y\log p_y$.

\begin{proposition}\label{prop:fisher}
For $u,v\in T$, $-d^2H_p(u,v)=\sum_yu_yv_y/p_y$, the Fisher metric.
\end{proposition}
\begin{proof}
$\partial^2H/\partial p_y\partial p_z=-\delta_{yz}/p_y$.
\end{proof}

The Fisher metric is symmetric, so it is not a symplectic form; the two structures coexist in the cotangent bundle, the
metric on the base and the canonical form $\omega=\sum dp_y\wedge d\xi_y$ on $T^*\Delta^\circ$, and on the toric manifolds of
\cref{sec:toric}, where the Fisher metric is twice the Kähler metric on the real section \cite[Prop.~6.5]{Conditionals}.

%=====================================================================
\section{Constraint branes in the cotangent bundle}\label{sec:conormal}
%=====================================================================

A marginal constraint fixes a linear image $Mp=b$ of the law. Its \emph{fibre} is $C=\{p\in\Delta:Mp=b\}$ and its open part is
$C^\circ=C\cap\Delta^\circ$. The fibre is not Lagrangian, but its conormal bundle is, as for every submanifold
\cite{McDuffSalamon2017}:
\[
N^*C^\circ=\{(p,\xi)\in T^*\Delta^\circ:p\in C^\circ,\ \xi\vert_{T_pC}=0\},\qquad T_pC=\ker M\cap T.
\]
The \emph{entropy brane} is the graph $L_H=\{(p,dH_p)\}$, an exact Lagrangian.

\begin{theorem}[Conormal intersection]\label{thm:conormal}
If $C^\circ\neq\emptyset$, then $L_H\cap N^*C^\circ$ is the single point over the maximum-entropy law $p^\star$ of $C$, and the
intersection is transversal. If $C^\circ=\emptyset$, then $N^*C^\circ=\emptyset$ and the intersection is empty, although $C$ itself may
be nonempty.
\end{theorem}
\begin{proof}
A point $(p,dH_p)$ lies in $N^*C^\circ$ iff $dH_p$ vanishes on $T_pC$, that is, iff $p$ is a critical point of $H$ on the
relatively open convex set $C^\circ$. $H$ is strictly concave, so there is at most one such point, and it is the maximizer of $H$
over $C$; the maximizer lies in $C^\circ$ when $C^\circ\ne\emptyset$, because the derivative of $H$ in a direction that raises a
vanishing coordinate is $+\infty$. Transversality of the graph of $dH$ to a conormal bundle at a point is nondegeneracy of the
Hessian of $H\vert_C$ there, and by \cref{prop:fisher} that Hessian is minus the Fisher metric on $T_pC$, which is definite.
\end{proof}

\begin{corollary}\label{cor:filler}
Suppose the margins admit a strictly positive joint law, so that $C^\circ\ne\emptyset$. For the chain margins $(X,Y)$ and $(Y,Z)$ the intersection point is the conditionally independent gluing
$P_{XY}P_{YZ}/P_Y$, the canonical filler of \cite[Thm.~4.2]{Conditionals}. For pairwise margins of three variables it is the
limit of iterative proportional fitting started at the uniform law \cite{Csiszar1975}.
\end{corollary}
\begin{proof}
The gluing has the margins, and $\log$ of it is a sum of a function of $(x,y)$ and one of $(y,z)$, which is orthogonal to the
kernel of the marginal map. The second statement is the maximum-entropy characterization of fitting.
\end{proof}

So the brane picture of compatibility is exact: a constraint is compatible with positive laws iff its brane is nonempty, and
then it meets the entropy brane exactly once, at the canonical filler.

\begin{example}[Three bits]\label{ex:threebits}
Let $X_0,X_1,X_2$ be bits, each uniform, and let each pair disagree with probability $1-\varepsilon$, so that each pairwise table is
$\bigl(\begin{smallmatrix}\varepsilon/2&(1-\varepsilon)/2\\(1-\varepsilon)/2&\varepsilon/2\end{smallmatrix}\bigr)$.
\begin{enumerate}[label=(\roman*)]
\item \emph{The fibre.} The kernel of the marginal map on $\R^8$ is spanned by the parity vector $x(s)=(-1)^{\lvert s\rvert}$, so the fibre
  is a segment along $x$.
\item \emph{Threshold.} Every assignment of three bits makes $0$ or $2$ of the three pairs disagree, so the expected number of
  disagreeing pairs, $3(1-\varepsilon)$, is at most $2$: a joint exists iff $\varepsilon\ge\frac13$. At $\varepsilon=\frac13$ the only joint puts no
  mass on $000$ and $111$, and at $\varepsilon=1$ the only joint lives on $000$ and $111$; in both cases $C^\circ=\emptyset$ and the brane is
  empty. A positive joint exists iff $\frac13<\varepsilon<1$. The two ends are the breaking regime of \cite{Bands}.
\item \emph{The filler.} For $\frac13<\varepsilon<1$ the maximum-entropy law gives $a=\frac{3\varepsilon-1}4$ to $000$ and $111$ and $b=\frac{1-\varepsilon}4$
  to the six other assignments, since then $\sum_sx(s)\log p(s)=0$.
\item \emph{Transversality and its loss.} The Fisher form of the fibre at the filler is
  $\sum_sx(s)^2/p(s)=\frac8{3\varepsilon-1}+\frac{24}{1-\varepsilon}$, which is $64$ at $\varepsilon=\frac12$, $197$ at $0.35$ and $436$ at $0.34$, and
  diverges at both ends: the intersection becomes infinitely steep before it disappears.
\item \emph{Contextual fraction.} Below the threshold the contextual fraction of \cite{AbramskyBarbosaMansfield2017} is $1-3\varepsilon$.
  A sub-normalized common joint $\lambda q$ below the tables puts at most $\varepsilon$ on the agreements of each pair, and every
  assignment agrees on some pair, so $\lambda\le3\varepsilon$; the measure $3\varepsilon$ times the uniform law on the six assignments that agree on
  exactly one pair attains it.
\end{enumerate}
\end{example}

\paragraph{Synthetic test: detecting the boundary.} Take the positive law with $\varepsilon=\frac13+m$ and $N$ joint samples of the three bits,
and estimate $\varepsilon$ by the mean agreement rate. Each sample agrees on one pair, or on three if it is $000$ or $111$, so $\hat\varepsilon\le\frac13$
exactly when no all-equal sample appears, which happens with probability $(1-\frac32m)^N\approx e^{-3Nm/2}$. For $m=0.02$ and $N=25,100,400$ the
verdict ``no positive filler'' was returned in $47\%$, $6\%$ and $0\%$ of $2000$ replicates, against $46.7\%$, $4.8\%$ and $0\%$ predicted. So
$O(1/m)$ joint samples suffice. The reason is that the filler's cells on $000$ and $111$ vanish at the boundary: the Fisher information of
$\varepsilon$ in the joint model, $\frac9{2(3\varepsilon-1)}+\frac3{2(1-\varepsilon)}$, diverges there, and the test reduces to seeing a rare cell.

The rate depends on the design. In the usual contextuality design each pair is observed in its own samples, the agreement count of each
pair is binomial, its Fisher information stays bounded at the boundary, and the misses decay only once $Nm^2$ is large: $42\%$, $26\%$, $7\%$
and $0.05\%$ for $N=25,100,400,1600$. Interior walls behave in the same way; for them \cite[Prop.~12.2]{Conditionals} shows that $N$ of order
$1/m^2$ suffices.

%=====================================================================
\section{Toric branes}\label{sec:toric}
%=====================================================================

For nondegenerate margins $f,g$, the coupling polytope $\Cpl(f,g)$ is Delzant for the lattice of integer circulations, and it is
the moment polytope of a compact toric manifold $X_{f,g}$ \cite[Thm.~6.4]{Conditionals}, \cite{Delzant1988}. The fibre of the moment
map over an interior coupling is a Lagrangian torus. A symplectic potential on the polytope determines a compatible toric
Kähler structure when it satisfies Abreu's conditions \cite{Abreu2003}: it differs from Guillemin's potential
$G_P=\frac12\sum_{\text{facets}}\ell\log\ell$ \cite{Guillemin1994} by a function smooth on the closed polytope, its Hessian is
positive definite inside, and $\det(\Hess G)\prod_{\text{facets}}\ell$ is smooth and positive on the closed polytope.

\begin{theorem}[Entropy is an admissible potential]\label{thm:abreu}
For nondegenerate margins, every cell $\alpha_c$ is either a facet function of $\Cpl(f,g)$ or strictly positive on it, and
$G=\frac12\sum_c\alpha_c\log\alpha_c=-\frac12H$ satisfies Abreu's conditions. Hence it defines a compatible toric Kähler structure on
$X_{f,g}$ whose metric on the real section is half the Fisher metric, and whose potential has the product coupling as its unique
critical point.
\end{theorem}
\begin{proof}
The polytope is simple, and at each vertex exactly $(n-1)(m-1)$ cells vanish, as many as the dimension. At a vertex of a simple
polytope given by a system in which exactly $d$ inequalities are tight, each tight inequality defines a facet. Every face on which
$\alpha_c=0$ contains a vertex where it is tight, so $\alpha_c$ is a facet function or never vanishes. In the second case
$\alpha_c\log\alpha_c$ is smooth on the polytope, so $G-G_P$ is smooth. The Hessian is $\frac12\sum_c u_c^2/\alpha_c$ on circulations, which is
positive definite. Near a vertex, in Delzant coordinates in which the facets through it are $x_i\ge0$ (each cell
function is primitive on the lattice of integer circulations, since a basic move takes the value $1$ on it, and distinct cells give
distinct facets off the walls), $\Hess G=\frac12X^{-1}+S$
with $X=\diag(x)$ and $S$ smooth and positive semidefinite, so $\det(\Hess G)\prod x_i=\det(\frac12I+XS)$ is smooth and at least
$2^{-d}$. A boundary point is covered by the coordinates of a vertex of its minimal face, in which the remaining facet functions
are positive factors. The critical point is computed in \cite[Prop.~6.5]{Conditionals}.
\end{proof}

This removes the hypothesis of \cite[Prop.~6.5]{Conditionals} that every cell be a facet. Guillemin's own potential is a
different admissible potential, and when some cells are not facets its critical point may differ from the product coupling
(\cref{ex:twobytwo}).

\begin{proposition}[The positive real Lagrangian]\label{prop:realpositive}
The positive real part $X_{f,g}^{>0}$, the fixed set of complex conjugation in the open torus orbit with positive coordinates, is a
Lagrangian submanifold, and the moment map restricts to a homeomorphism of its closure onto $\Cpl(f,g)$ \cite{Fulton1993}. It meets the
torus brane over each interior coupling in exactly one point. So the couplings are the points of one Lagrangian, and also the
intersection points of that Lagrangian with the torus branes.
\end{proposition}
\begin{proof}
Complex conjugation is antiholomorphic and preserves the torus-invariant Kähler form up to sign, so its fixed set is Lagrangian. The
homeomorphism is the classical statement about the nonnegative part of a projective toric variety; a torus fibre meets the positive
part in the single point whose angle coordinates are all zero.
\end{proof}

\begin{proposition}[Which torus branes survive]\label{prop:floer}
Suppose $X_{f,g}$ is Fano. The torus brane over a coupling $u$, with a unitary local system, has nonzero Floer cohomology iff $u$ and the local system give a critical
point of the disc potential $\mathfrak{PO}(u;y)=\sum_{\text{facets}}T^{\ell(u)}y^{\nu}$, the Cho--Oh formula for toric Fano manifolds
\cite{ChoOh2006,FOOO2010}. All other torus branes are zero objects of the Fukaya category.
\end{proposition}

\begin{example}[$2\times2$ couplings]\label{ex:twobytwo}
Let $f=(0.4,0.6)$ and $g=(0.3,0.7)$. With $t=\alpha_{11}$ the polytope is the segment $[0,0.3]$, the cells are
$(t,\ 0.4-t,\ 0.3-t,\ 0.3+t)$, and $X_{f,g}=\mathbb{CP}^1$.
\begin{enumerate}[label=(\roman*)]
\item The facets are $t\ge0$ and $0.3-t\ge0$; the cells $0.4-t$ and $0.3+t$ stay at least $0.1$ and $0.3$, as \cref{thm:abreu} predicts.
\item The entropy potential is critical at $t=0.12$, the product coupling $f_1g_1$; Guillemin's potential is critical at the midpoint
  $t=0.15$. The two Kähler structures differ, although their symplectic forms agree.
\item $\mathfrak{PO}=T^{t}y+T^{0.3-t}y^{-1}$ has a critical point with unitary $y$ only when both terms have equal valuation, $t=0.15$,
  with $y=\pm1$. So the only
  nonzero torus brane is over the Fréchet midpoint $\frac12(\max(0,f_1+g_1-1)+\min(f_1,g_1))$, the circle that divides the sphere into
  equal areas. Here the brane over the product coupling is a zero object, as in \cite[Prop.~10.11]{Conditionals}. The two can
  coincide: for $f=(0.4,0.6)$ and $g=(0.5,0.5)$ the segment is $[0,0.4]$ and the product coupling $0.2$ is also the midpoint and
  Guillemin's critical point.
\end{enumerate}
\end{example}

%=====================================================================
\section{Correspondences: where gluing stops being geometric}\label{sec:correspondences}
%=====================================================================

Gluing with a fixed $\beta\in\Cpl(g,h)$ is the map $\alpha\mapsto\alpha\,\diag(g)^{-1}\beta=\alpha K$, where $K=\diag(g)^{-1}\beta$ is the
kernel of $\beta$. It is linear, and it carries $\Cpl(f,g)$ into $\Cpl(f,h)$ with $h=K^\T g$. The tori of $X_{f,g}$ and $X_{f,h}$ have the
integer circulations of the two shapes as weight lattices \cite[Thm.~6.4]{Conditionals}. Call a Lagrangian correspondence
$\Gamma\subset X_{f,g}^-\times X_{f,h}$ \emph{equivariant} if it is connected and invariant under the torus $T_{f,h}$, acting on $X_{f,h}$ and on
$X_{f,g}$ through a homomorphism $\psi\colon T_{f,h}\to T_{f,g}$, and say that it \emph{induces} a map $\Phi$ of polytopes if
$(\mu(z),\mu(z'))=(\alpha,\Phi(\alpha))$ for points $(z,z')\in\Gamma$ over a full-dimensional set of $\alpha$. The moment map of the diagonal
action is constant on the connected invariant Lagrangian $\Gamma$, so $\mu_{f,h}=d\psi^*\mu_{f,g}+c$ on $\Gamma$, and $\Phi$ has linear part
$d\psi^*$, which maps weight lattice to weight lattice.

\begin{theorem}[Integrality]\label{thm:integral}
The linear map $x\mapsto xK$ carries integer $n\times m$ circulations to integer $n\times h$ circulations iff $K$ is deterministic,
every row a vertex of the simplex, or constant, all rows equal. So gluing with a coupling whose kernel is stochastic and not constant
is not induced by any equivariant Lagrangian correspondence in the sense above.
\end{theorem}
\begin{proof}
Integer circulations are generated by the moves $(e_i-e_k)(e_j-e_l)^\T$, and such a move goes to $(e_i-e_k)(K_j-K_l)$. So the condition is
that any two rows of $K$ differ by an integer vector. Two stochastic rows with integer difference are equal or are distinct vertices;
if some row is not a vertex, every other row equals it.
\end{proof}

\begin{example}[$2\times2$ gluing]\label{ex:gluetwo}
For $2\times2$ couplings the circulations form the line spanned by $x=\bigl(\begin{smallmatrix}1&-1\\-1&1\end{smallmatrix}\bigr)$, and
$xK=\lambda x$ with $\lambda=K_{11}-K_{21}$. Its absolute value is the Dobrushin coefficient of $K$. The moment segment is carried onto a
segment $\lvert\lambda\rvert$ times as long, so the map of moment polytopes contracts symplectic area by $\lvert\lambda\rvert$, and it is integral
exactly for $\lambda\in\{1,-1,0\}$: the identity, the swap and a constant kernel. In \cref{ex:twobytwo}, a noisy relabelling with
$K=\bigl(\begin{smallmatrix}0.9&0.1\\0.2&0.8\end{smallmatrix}\bigr)$ has $\lambda=0.7$ and shrinks the segment from length $0.3$ to $0.21$.
\end{example}

\begin{remark}[Related defects]\label{rem:split}
Stochastic kernels cause related, though not identical, defects elsewhere in the programme: whiskering is functorial for
deterministic kernels and fails for stochastic ones in general \cite[Thm.~10.2]{Conditionals}; squares of couplings and kernels
compose for bijections and sufficient merges, and fail for some non-injective functions and some constant kernels
\cite[Prop.~6.1]{Algebras}; and kernels contract bands by their Dobrushin coefficient \cite{Bands}. The class singled out differs from
case to case; here constant kernels are on the good side, there they are not. What the cases share is that the defect is a
circulation, or a contraction of the circulation space.
\end{remark}

%=====================================================================
\section{The categorical layer}\label{sec:categorical}
%=====================================================================

Couplings can enter a categorical layer as morphisms or as objects, and each choice loses something.

\begin{proposition}[Morita data]\label{prop:morita}
The gluing algebra $A_f\cong\R\times M_{n-1}(\R)$ \cite[Thm.~2.2]{Algebras} is Morita equivalent to $\R\times\R$. Hence its
category of perfect complexes is that of two points, $K_0(A_f)=\Z^2$, its Hochschild cohomology is $\R^2$ in degree $0$ (its centre, spanned
by $P_f$ and $e_f$) and zero above, and its simple modules have dimensions $1$ and $n-1$.
\end{proposition}
\begin{proof}
$M_{n-1}(\R)$ is Morita equivalent to $\R$, and Morita equivalence preserves derived categories, $K$-theory and Hochschild
cohomology \cite{Keller2006}. The centre of a product of matrix algebras is the product of the centres.
\end{proof}

\begin{center}
\begin{tabular}{p{3.2cm}p{4.6cm}p{6.2cm}}
\toprule
couplings as & composition & what survives\\
\midrule
morphisms (spans) & exactly gluing \cite{Algebras} & for functors that kill $e_f$ up to homotopy, only the mass \cite[Cor.~5.3]{Algebras}; derived invariants of two points\\
torus branes & not equivariantly induced for non-constant stochastic kernels (\cref{thm:integral}) & in the Fano case, only critical points of $\mathfrak{PO}$ (\cref{prop:floer})\\
points of $X^{>0}_{f,g}$ & none & all couplings, with positivity (\cref{prop:realpositive})\\
\bottomrule
\end{tabular}
\end{center}

Positivity, where contextuality lives \cite[Cor.~5.3]{Algebras}, is invisible in the first two rows: linear spans forget it, and
which torus branes survive is decided by the disc potential, not by positivity. It is kept in the third row. So the promising
structure is the Floer theory between the positive real Lagrangian and the torus branes.

%=====================================================================
\section{Open problems}\label{sec:open}
%=====================================================================

\begin{enumerate}[leftmargin=*]
\item \emph{Real and torus branes.} Compute $HF(X^{>0}_{f,g},L_u)$, or of the full real locus, with its products, and decide whether the
  positivity of couplings survives in it.
\item \emph{Non-equivariant correspondences.} \Cref{thm:integral} rules out equivariant ones. Is there a Lagrangian correspondence, not
  torus-equivariant, whose composition matches gluing with a stochastic kernel?
\item \emph{Bulk deformations.} With bulk deformation the set of nonzero torus branes changes. Is there a bulk class for which the
  product coupling's brane is nonzero for all margins?
\item \emph{Larger tables.} Locate the nonzero torus branes for $n\times m$ couplings, extending \cref{ex:twobytwo}; the computations of
  \cite[Obs.~10.10]{Conditionals} point to the leximin coupling.
\item \emph{The three-bit brane under noise.} Track the conormal brane of \cref{ex:threebits} as a family over $\varepsilon$, through the
  threshold, as a first example of a brane that disappears at a breaking point.
\end{enumerate}

%=====================================================================
\section{Computational checks}\label{sec:code}
%=====================================================================

\texttt{check\_bridges.py} runs nine checks, all passing: G1 (\cref{prop:fisher}); G2 (\cref{thm:conormal,cor:filler}); W1
(\cref{ex:threebits}, all five parts); S1 (the synthetic detection test, with joint and with pairwise-only samples); G3 (\cref{thm:abreu}, the determinant condition at every vertex
of an example with non-facet cells); W2 (\cref{ex:twobytwo}); G4 (\cref{thm:integral}); W3 (\cref{ex:gluetwo}); W4 (\cref{prop:morita}, the
dimension and the centre).

\begin{thebibliography}{99}

\bibitem{Algebras}
V.~K. Awasthi, Coupling algebras: gluing, the square-zero cone and the ring of contingency tables. Technical report, Johns
Hopkins University, October 2026. Draft 4.

\bibitem{Bands}
V.~K. Awasthi, Bands: perturbation, tearing and breaking of coupling structures, and the entropy of changing probabilities.
Technical report, Johns Hopkins University, October 2026. Draft 1.

\bibitem{Conditionals}
V.~K. Awasthi, Conditionals, couplings and toric geometry: kernels, the coupling nerve, cumulants and a {F}ukaya
programme. Technical report, Johns Hopkins University, October 2026. Revised draft 4.

\bibitem{AbramskyBarbosaMansfield2017}
S.~Abramsky, R.~S. Barbosa and S.~Mansfield, Contextual fraction as a measure of contextuality. \emph{Phys. Rev. Lett.}
\textbf{119} (2017), 050504.

\bibitem{Abreu2003}
M.~Abreu, Kähler geometry of toric manifolds in symplectic coordinates. In \emph{Symplectic and Contact Topology}, Fields Inst.
Commun. 35, Amer. Math. Soc., 2003, 1--24.

\bibitem{ChoOh2006}
C.-H. Cho and Y.-G. Oh, Floer cohomology and disc instantons of Lagrangian torus fibers in Fano toric manifolds. \emph{Asian J.
Math.} \textbf{10} (2006), 773--814.

\bibitem{Csiszar1975}
I.~Csiszár, {$I$}-divergence geometry of probability distributions and minimization problems. \emph{Ann. Probab.} \textbf{3}
(1975), 146--158.

\bibitem{Delzant1988}
T.~Delzant, Hamiltoniens périodiques et images convexes de l'application moment. \emph{Bull. Soc. Math. France} \textbf{116}
(1988), 315--339.

\bibitem{FOOO2010}
K.~Fukaya, Y.-G. Oh, H.~Ohta and K.~Ono, Lagrangian Floer theory on compact toric manifolds {I}. \emph{Duke Math. J.} \textbf{151}
(2010), 23--174.

\bibitem{Fulton1993}
W.~Fulton, \emph{Introduction to Toric Varieties}, Annals of Mathematics Studies 131, Princeton University Press, 1993.

\bibitem{Guillemin1994}
V.~Guillemin, Kaehler structures on toric varieties. \emph{J. Differential Geom.} \textbf{40} (1994), 285--309.

\bibitem{Keller2006}
B.~Keller, On differential graded categories. In \emph{Proceedings of the International Congress of Mathematicians, Madrid 2006},
Vol.~II, European Mathematical Society, 2006, 151--190.

\bibitem{McDuffSalamon2017}
D.~McDuff and D.~Salamon, \emph{Introduction to Symplectic Topology}, 3rd ed., Oxford University Press, 2017.

\end{thebibliography}
\end{document}
